\section{Quantitative Analysis}

\subsection{Calibration Strategy}

This section maps the theoretical model to a quantitative framework and describes the calibration strategy. The objective is not to match a specific country exactly, but to choose parameters that generate empirically plausible magnitudes and clearly illustrate the model's mechanisms.

The calibration focuses on economies with fixed or tightly managed exchange rate regimes where parallel markets are active.

\subsection{Mapping from Model to Quantitative Objects}

The key objects in the model map to observable or interpretable counterparts as follows:

\begin{itemize}
\item Fundamental $\theta$: deviation of the official exchange rate from the shadow (market-clearing) rate
\item Private signal $x_i$: noisy cash-based parallel market price
\item Public signal $s$: stablecoin price in local currency terms
\item Stablecoin usage $q_s$: trading volume or adoption intensity
\item Crisis probability $\pi$: probability of a discrete devaluation or regime break
\end{itemize}

\subsection{Parameters}

We divide parameters into three groups.

\paragraph{Information parameters}
\begin{align}
\sigma_x^2 &\text{: noise in cash-based signal} \\
\bar{\sigma}_s^2 &\text{: baseline noise in stablecoin signal} \\
\rho &\text{: sensitivity of signal precision to stablecoin usage}
\end{align}

\paragraph{Preference and payoff parameters}
\begin{align}
\zeta &\text{: utility from foreign currency holdings} \\
\lambda &\text{: loss severity in crisis} \\
\xi &\text{: amplification from stablecoin exposure}
\end{align}

\paragraph{Policy parameters}
\begin{align}
\alpha &\text{: strength of reserves or policy buffer} \\
\delta &\text{: sensitivity of regime stability to misalignment}
\end{align}

\subsection{Calibration Targets}

Parameters are chosen to match the following stylized facts:

\begin{enumerate}
\item Stablecoin usage increases with parallel market premia
\item Crisis probabilities are low but non-zero in normal times
\item Welfare gains from financial innovation are moderate, not extreme
\end{enumerate}

\subsection{Solution Method}

We solve the model using two approaches.

\paragraph{Full solver}

The full solver computes a fixed point in which:
\begin{itemize}
\item asset demand determines signal precision
\item signal precision determines beliefs
\item beliefs determine aggregate behavior and crisis probability
\end{itemize}

This corresponds exactly to the theoretical model.

\paragraph{Fast solver}

The fast solver approximates the equilibrium by simplifying the mapping between signal precision and behavior while preserving the key comparative statics.

The fast solver is used for calibration and sensitivity analysis, while the full solver is used to validate the results.

\subsection{Consistency with Theory}

All quantitative results are consistent with the theoretical model. In particular:
\begin{itemize}
\item Signal precision is endogenous and determined by $q_s$
\item Crisis probability depends on aggregate behavior
\item Welfare reflects both efficiency and fragility components
\end{itemize}

Thus, the quantitative exercise is a direct implementation of the theoretical framework rather than a reduced-form approximation.
